% Part: first-order-logic
% Chapter: beyond
% Section: other-logics

\documentclass[../../../include/open-logic-section]{subfiles}

\begin{document}

\olfileid{fol}{byd}{oth}

\olsection{इतर तर्कशास्त्रे}

आतापर्यंत तुमच्या लक्षात आले असेलच की नवे तर्कशास्त्र रचणे अवघड नाही. तुम्हीही
स्वतःचा विन्यास तयार करू शकता, निगमन-पद्धती बनवू शकता आणि तिला साजेशी
चिन्हार्थमीमांसा घडवू शकता. !!{derivation} प्रणाली चिन्हार्थमीमांसेसाठी
संपूर्ण हवी असेल, तर थोडे चातुर्य वापरावे लागू शकते; तसेच तुमचे तर्कशास्त्र
खरोखर रंजक आहे हे जगाला पटवून देण्यासाठी काही प्रयत्न करावे लागू शकतात. पण
त्याबदल्यात त्याच्या गणितीय आणि संगणकीय गुणधर्मांचा शोध घेताना अनेक तास
निर्भेळ आनंद मिळवता येईल.

अलीकडच्या दशकांत आकारिक तर्कशास्त्रांचा खरोखरच स्फोट झाला आहे. संदिग्ध
गुणधर्मांविषयीच्या तर्कविचाराचे प्रतिमान देण्यासाठी अस्पष्ट तर्कशास्त्र रचले
आहे. अनिश्चिततेविषयीच्या तर्कविचाराचे प्रतिमान देण्यासाठी संभाव्यताधारित
तर्कशास्त्र रचले आहे. खंडनीय प्रकारच्या तर्कविचाराचे प्रतिमान देण्यासाठी
पूर्वमान्य तर्कशास्त्रे आणि अ-एकस्वनिक तर्कशास्त्रे रचली आहेत; म्हणजे नवी
माहिती मिळाल्यावर नंतर उलटवता येतील असे “वाजवी” निष्कर्ष. ज्ञानाविषयीच्या
तर्कविचाराचे प्रतिमान देणारी ज्ञानविषयक तर्कशास्त्रे आहेत; कार्यकारण
संबंधांविषयीच्या तर्कविचाराचे प्रतिमान देणारी कार्यकारणविषयक तर्कशास्त्रे
आहेत; आणि नैतिक व नीतिविषयक कर्तव्यांविषयीच्या तर्कविचाराचे प्रतिमान देणारी
“कर्तव्यविषयक” तर्कशास्त्रेही आहेत. या प्रणाली मांडण्यामागील मुख्य प्रेरणा
तात्त्विक, गणितीय किंवा संगणकीय आहे यानुसार, अशा गोष्टींचा अभ्यास गणितीय
तर्कशास्त्र, तात्त्विक तर्कशास्त्र, कृत्रिम बुद्धिमत्ता, संज्ञानविज्ञान किंवा
इतर क्षेत्रांच्या नावाखाली झालेला आढळेल.

ही यादी सतत वाढत जाते आणि शक्यता अंतहीन दिसतात. सर्व मानवी तर्कबुद्धीचे गणनेत
रूपांतर करण्याचे लायप्निट्स यांचे स्वप्न आपण कदाचित कधीच साध्य करणार नाही---
पण त्यामुळे प्रयत्न करणे थांबवण्याचे कारण नाही.

\end{document}
